\documentclass[11pt,a4paper]{article}
\usepackage[utf8]{inputenc}
\usepackage[italic]{mathastext}
\usepackage[italian]{babel}
\usepackage{geometry}
\geometry{a4paper, margin=20mm}
\usepackage{hyperref}
\usepackage{xcolor}
\usepackage{booktabs}
\usepackage{listings}
\usepackage{enumitem}

\definecolor{primary}{RGB}{15, 23, 42}
\definecolor{accent}{RGB}{217, 119, 6}
\definecolor{codebg}{RGB}{15, 23, 42}

\hypersetup{
    colorlinks=true,
    linkcolor=primary,
    urlcolor=accent
}

\lstset{
    backgroundcolor=\color{codebg},
    basicstyle=\ttfamily\small\color{white},
    breaklines=true,
    frame=single,
    rulecolor=\color{accent},
    language=SQL
}

\title{\textbf{\Large RELAZIONE TECNICA DETTAGLIATA DI PROGETTO}\\ \large Corso: Web and Multimedia Technologies — UNIPV}
\author{\textbf{Marco Santoro}}
\date{Anno Accademico 2025/2026}

\begin{document}

\maketitle

\section{Introduzione, Vision e Contesto Applicativo}
La presente relazione descrive in dettaglio la progettazione, l'architettura tecnica e le scelte metodologiche alla base dello sviluppo del sito web e della piattaforma di gestione per \textbf{Algora Studio}. Algora Studio è una software house specializzata nella realizzazione di soluzioni verticali per il dominio del patrimonio culturale e della memoria storica documentale italiana.

Il prodotto di punta è \textbf{Foliarium}, un software gestionale sviluppato per l'Archivio di Stato di Savona per la digitalizzazione e consultazione \textit{fuzzy} degli archivi catastali storici (dal 1830 ad oggi).

\section{Architettura dell'Informazione e Mappatura delle Pagine}
L'applicazione si articola su 5 pagine HTML5 statiche integrate da 1 script dinamico PHP:

\begin{table}[h!]
\centering
\small
\begin{tabular}{lll}
\toprule
\textbf{Risorsa File} & \textbf{Stack} & \textbf{Ruolo nell'Architettura} \\
\midrule
\texttt{index.html} & HTML5, CSS3, JS & Home Page Istituzionale con manifesto e metriche. \\
\texttt{chi-siamo.html} & HTML5, CSS3 & Profilo aziendale e filosofia dello sviluppo verticale. \\
\texttt{foliarium.html} & HTML5, CSS3 & Scheda prodotto dettagliata e prezzi licenze PA. \\
\texttt{caso-studio.html} & HTML5, CSS3 & Caso studio reale Archivio di Stato di Savona. \\
\texttt{contatti.html} & HTML5, CSS3, JS & Interfaccia di contatto, FAQ accordion e form. \\
\texttt{invia-contatto.php} & PHP 8, MySQLi & Backend server-side per la persistenza dei dati. \\
\bottomrule
\end{tabular}
\end{table}

\section{Tecniche Frontend: HTML5, CSS3 e Layout Non Lineare}
Il codice è stato sviluppato interamente a mano. Il layout "non lineare" è garantito dall'uso di \texttt{position: fixed} per la navigazione e da griglie asimmetriche in \textbf{CSS Grid} e \textbf{Flexbox}.

\section{Tecniche Client-Side: JavaScript Vanilla}
Realizzato in JS Vanilla per la gestione del menu responsive (\texttt{toggleNav}), dello scroll listener e del sistema FAQ accordion.

\section{Tecniche Server-Side e Persistenza Dati: PHP \& MySQL}
Schema DDL della tabella \texttt{contatti}:

\begin{lstlisting}
CREATE DATABASE IF NOT EXISTS algora_db;
USE algora_db;

CREATE TABLE IF NOT EXISTS contatti (
    id INT AUTO_INCREMENT PRIMARY KEY,
    nome VARCHAR(100) NOT NULL,
    ente VARCHAR(100),
    email VARCHAR(100) NOT NULL,
    tipo VARCHAR(50),
    messaggio TEXT NOT NULL,
    data_invio DATETIME DEFAULT CURRENT_TIMESTAMP
) ENGINE=InnoDB;
\end{lstlisting}

Sicurezza garantita da Prepared Statements (\texttt{mysqli}) e sanitizzazione degli input (\texttt{trim()}, \texttt{htmlspecialchars()}).

\section{Conclusioni e Collaudo Locale}
Per il collaudo locale: avviare Apache e MySQL su XAMPP, copiare la cartella in \texttt{htdocs}, importare il database e accedere a \texttt{http://localhost/algora_site/contatti.html}.

\end{document}
